% problem-000116 5 元素推断与化合物
\section*{5. 元素推断与化合物}
⾦属氮化物材料具有⾼硬度、⾼熔点、⾼热稳定性特点。最近，科学家研究的某三元氮化物表现出特殊导电性质。该晶体结构如图4所⽰。每⼀层氮原⼦都是密置层，且层间距相等。铁原⼦和钨原⼦相间填充在两层氮原⼦之间，钨原⼦填充在氮原⼦形成的三棱柱配位体中⼼，铁原⼦填充在氮原⼦形成的⼋⾯配位体中⼼。

（图：）
图4 三元氮化物结构

\subsection*{5-1}
请根据晶体结构画出⼀个晶胞（N放在顶点位置，以●表⽰N原⼦，以◯表⽰Fe原⼦，以ⓧ表⽰W原⼦），并写出晶胞中⾦属原⼦的分数坐标。

\subsection*{5-2}
分别⽤ $\mathrm { A / B / C }$ 表⽰氮原⼦，a/b/c表⽰铁原⼦， $\mathrm { a ^ { \prime } / b ^ { \prime } / c ^ { \prime } }$ 表⽰钨原⼦的堆积位置，请表⽰出沿�轴⽅向原⼦的堆积⽅式（如 $\mathrm { { A c B a ^ { \prime } C b A } ) }$ o

\subsection*{5-3}
已知晶胞参数 $a \ = 0 . 2 8 7 6 \ \mathrm { n m } , \ c \ = 1 . 0 9 3 \ \mathrm { n m } \circ$

\subsection*{5-3}
-1 计算该晶体的密度�。

\subsection*{5-3}
-2 计算该晶体中W原⼦与N原⼦的最短核间距�。

\subsection*{5-4}
该晶体中⾦属原⼦与氮原⼦之间的作⽤⼒有明显的共价成分，表现出各向异性的⾦属导电性，请指出晶体的导电⽅向，并简要说明理由。
